\begin{abstract}
\noindent
Primal--dual active-set methods and block principal pivoting solve a strictly
convex quadratic program by guessing a whole active set at once and repairing it
from the signs the working-set subproblem returns. For bound constraints the
guess is safe: the system solved on a candidate set is a principal submatrix of
$G^{-1}$, positive definite whatever the guess, and that $P$-matrix property is
the hypothesis under which the exchange terminates finitely. For general
constraints $C\T x \ge b$ the hypothesis fails. The system is
$C_\Active\T G^{-1} C_\Active$, positive definite exactly when $C_\Active$ has
full column rank, which is a property of the guess and not of the data.

That the plain exchange needs repairing on general constraints is known: Curtis,
Han and Robinson recover global convergence by carrying an index set auxiliary to
the active-set estimate. What we add is the mechanism and its price. The
obstruction is structural rather than a missing anti-cycling rule, because no rule
for choosing the next set can recover the bound-constrained theorem when the
subproblem on the chosen set may be singular; and the singular guess is reached by
the method's own starting rule, not only by an adversarial one. What keeps such a
method usable is that strict convexity makes the Karush--Kuhn--Tucker conditions
sufficient, so a candidate set is certified rather than trusted.

The exposure is easy to mistake for a technicality, so we measure it on a shipped
solver. Over the $600$ instances of four constraint families it never
occurs, which is why it can be overlooked. Duplicating columns of $C$, as any
programmatically assembled model may, raises it to $\qpPdasSingularDup\%$, and the
certified fraction on that family falls from $\qpPdasCertSmall\%$ at
$n = \qpPdasNsmall{}$ to $\qpPdasCertBig\%$ at $n = \qpPdasNbig{}$, against
$100\%$ everywhere else.
\end{abstract}
